\documentclass{article}
\usepackage{lmodern}
\usepackage{amsmath}
\usepackage{listings}
\usepackage{fullpage}
\usepackage{parskip}
\usepackage{xcolor}
\lstset{
	basicstyle=\ttfamily,
	columns=fullflexible,
	frame=single,
	breaklines=true,
	postbreak=\mbox{\textcolor{red}{$\hookrightarrow$}\space},
}
\begin{document}
\title{Experiment `p\_4\_4\_month\_with\_highest\_temperature\_array' Results}
\author{\today}
\date{}
\maketitle
\textbf{Experiment outcome: }SUCCESS
\\\textbf{Bad responses: }0
\\\textbf{Responses containing}~\texttt{assume}~\textbf{: }0
\\\textbf{Responses containing}~\texttt{expect}~\textbf{: }0
\\\textbf{Resolution attempts: }1
\\\textbf{Hard fails (resolution): }0
\\\textbf{Soft fails (resolution): }0
\\\textbf{Verification attempts: }1
\section*{Problem Specification}
\textbf{Problem name: }p\_4\_4\_month\_with\_highest\_temperature\_array
\\\textbf{Natural language statement: }Write a method that reads a sequence of twelve temperature values and returns the number of the month with the highest temperature.
\\\textbf{Method signature: }p\_4\_4\_month\_with\_highest\_temperature\_array(inputs: array<real>) returns (hottest: int)
\subsection*{Ensures}
\begin{itemize}
\item \texttt{1 <= hottest <= 12}
\item \texttt{forall i :: 0 <= i <= inputs.Length - 1 ==> inputs[hottest - 1] >= inputs[i]}
\end{itemize}
\subsection*{Requires}
\begin{itemize}
\item \texttt{inputs.Length == 12}
\end{itemize}
\clearpage
\section*{GenAI interactions}
Below you will find all interactions between the `user' (program) and the `assistant' (OpenAI).
\subsection*{Program $\rightarrow$ GenAI}
\begin{lstlisting}
You are given the following task to perform in Dafny:

Write a method that reads a sequence of twelve temperature values and returns the number of the month with the highest temperature.

The signature should be:

p_4_4_month_with_highest_temperature_array(inputs: array<real>) returns (hottest: int)

The method should respect the following contract:

requires inputs.Length == 12, ensures 1 <= hottest <= 12, ensures forall i :: 0 <= i <= inputs.Length - 1 ==> inputs[hottest - 1] >= inputs[i]

Produce and show only the Dafny body of this method, including the curly braces that surround it. Do not show the signature nor contract. You must not use 'assume' nor 'expect' anywhere in your code.
\end{lstlisting}
\subsection*{GenAI $\rightarrow$ Program}
\textbf{System fingerprint: }null
\\\textbf{ID: }chatcmpl-DmfXBEkQvr0YD0tgd5fnCW5eVkt2k
\\\textbf{Model: }gpt-5.5-2026-04-23
\\\textbf{Created at: }1780491773
\\\textbf{Finish reason: }stop
\\\textbf{Usage: }354 tokens in, and 646 tokens out
\begin{lstlisting}
{
  hottest := 1;
  var idx := 1;
  while idx < inputs.Length
    invariant 1 <= idx <= inputs.Length
    invariant 1 <= hottest <= idx
    invariant forall i :: 0 <= i < idx ==> inputs[hottest - 1] >= inputs[i]
    decreases inputs.Length - idx
  {
    if inputs[idx] > inputs[hottest - 1] {
      hottest := idx + 1;
    }
    idx := idx + 1;
  }
}
\end{lstlisting}
\section*{Final Program}
\begin{lstlisting}
method p_4_4_month_with_highest_temperature_array(inputs: array<real>) returns (hottest: int)
	requires inputs.Length == 12
	ensures 1 <= hottest <= 12
	ensures forall i :: 0 <= i <= inputs.Length - 1 ==> inputs[hottest - 1] >= inputs[i]
{
  hottest := 1;
  var idx := 1;
  while idx < inputs.Length
    invariant 1 <= idx <= inputs.Length
    invariant 1 <= hottest <= idx
    invariant forall i :: 0 <= i < idx ==> inputs[hottest - 1] >= inputs[i]
    decreases inputs.Length - idx
  {
    if inputs[idx] > inputs[hottest - 1] {
      hottest := idx + 1;
    }
    idx := idx + 1;
  }
}
\end{lstlisting}
\section*{Total Token Usage}
\textbf{Input tokens: }354
\\\textbf{Output tokens: }646
\\\textbf{Reasoning tokens: }512
\\\textbf{Sum of `total tokens': }1000
\section*{Experiment Timings}
\textbf{Overall Experiment} started at 1780491773517, ended at 1780491793124, lasting 19607ms (19.61 seconds)
\\\textbf{Iteration \#1} started at 1780491773517, ended at 1780491793124, lasting 19607ms (19.61 seconds)
\end{document}
